% =====================================================================
%  Theory: decomposition, identification, sample complexity, conformal
% =====================================================================
\section{Theory}
\label{sec:theory}

\subsection{Setup and notation}

Let $P_s$ and $P_t$ be distributions on $\mathcal{X}\times\{0,1\}$, where
$\mathcal{X}$ is the space of 12-lead recordings and $Y$ indicates one
harmonised label (labels are treated one-versus-rest throughout; the multiclass
statements are given in Appendix~\ref{app:multiclass}).  Write
$\pi_s=P_s(Y=1)$ and $\pi_t=P_t(Y=1)$.

Let $f:\mathcal{X}\to[0,1]$ be the frozen source-trained scorer, and let
$\mu_t$ denote the law of $V=f(X)$ under $P_t$.  Define the target
\emph{calibration curve}
\[
  g_t(v)\;=\;\mathbb{E}_{P_t}\!\left[\,Y \mid f(X)=v\,\right],
\]
and the calibration errors
\[
  \mathrm{CE}_1(f;P_t)=\big\|g_t-\mathrm{id}\big\|_{L^1(\mu_t)},
  \qquad
  \mathrm{CE}_2^2(f;P_t)=\big\|g_t-\mathrm{id}\big\|^2_{L^2(\mu_t)} .
\]
$\mathrm{CE}_1$ is the population quantity the empirical ECE estimates;
$\mathrm{CE}_2^2$ is the reliability term of the Brier decomposition.

\begin{assumption}[Source calibration]\label{ass:srccal}
$f(x)=P_s(Y=1\mid X=x)$ for $P_s$-almost every $x$.
\end{assumption}

Assumption~\ref{ass:srccal} is not cosmetic and not free: it is enforced in the
pipeline by temperature scaling on a held-out source split before the model is
frozen.  Without it, residual source miscalibration would be carried into every
target site and attributed by the decomposition below to shift.

\begin{definition}[Prior-correction map]\label{def:T}
For $w_1=\pi_t/\pi_s$ and $w_0=(1-\pi_t)/(1-\pi_s)$, let
\[
  T(v)\;=\;T_{\pi_s\to\pi_t}(v)\;=\;\frac{w_1 v}{w_1 v + w_0 (1-v)} .
\]
\end{definition}

$T$ depends on the data only through the two priors; in particular it requires
\emph{no target labels}, since $\pi_t$ is estimable from unlabeled target
scores (\S\ref{sec:ident}).

\subsection{Theorem 1: decomposition}

\begin{definition}[Shift components]
$\delta_{\mathrm{lab}} = T-\mathrm{id}$ and
$\delta_{\mathrm{con}} = g_t-T$, both as elements of $L^2(\mu_t)$.
\end{definition}

\begin{theorem}[Decomposition]\label{thm:decomp}
Under Assumption~\ref{ass:srccal},
\begin{enumerate}[label=(\alph*)]
\item \textbf{(Exact $L^2$ identity.)}
\[
  \mathrm{CE}_2^2(f;P_t)
  \;=\;
  \underbrace{\|\delta_{\mathrm{lab}}\|^2_{L^2(\mu_t)}}_{D_{\mathrm{label}}}
  \;+\;
  \underbrace{\|\delta_{\mathrm{con}}\|^2_{L^2(\mu_t)}}_{D_{\mathrm{concept}}}
  \;+\;
  \underbrace{2\,\langle \delta_{\mathrm{lab}},\delta_{\mathrm{con}}\rangle_{L^2(\mu_t)}}_{R} .
\]
\item \textbf{($L^1$ bound and its tightness.)}
$\mathrm{CE}_1(f;P_t)\le \|\delta_{\mathrm{lab}}\|_{L^1}+\|\delta_{\mathrm{con}}\|_{L^1}$,
with equality if and only if
$\operatorname{sign}\delta_{\mathrm{lab}}=\operatorname{sign}\delta_{\mathrm{con}}$
$\mu_t$-almost everywhere.
\item \textbf{(The concept term is exactly the failure of label shift on
$\sigma(f)$.)}  $\delta_{\mathrm{con}}=0$ $\mu_t$-a.e.\ if and only if
$P_t(Y=1\mid f(X))=T(f(X))$ a.s.  In particular, if
$P_t(x\mid y)=P_s(x\mid y)$ for $y\in\{0,1\}$ then $\delta_{\mathrm{con}}=0$,
so that $\mathrm{CE}_2^2 = D_{\mathrm{label}}$ and the whole miscalibration is
removed by $T$.
\end{enumerate}
\end{theorem}

\begin{proof}
(a) $g_t-\mathrm{id}=\delta_{\mathrm{lab}}+\delta_{\mathrm{con}}$ by
construction, and the identity is the expansion of
$\|a+b\|^2=\|a\|^2+\|b\|^2+2\langle a,b\rangle$ in $L^2(\mu_t)$.

(b) The triangle inequality in $L^1$, with the standard equality condition for
$|a+b|=|a|+|b|$ pointwise.

(c) The ``only if'' direction is the definition.  For the ``if'': suppose
$P_t(x\mid y)=P_s(x\mid y)$.  Then for $P_t$-a.e.\ $x$,
\[
  \frac{P_t(Y=1\mid X=x)}{P_t(Y=0\mid X=x)}
  =\frac{\pi_t\,p_s(x\mid 1)}{(1-\pi_t)\,p_s(x\mid 0)}
  =\frac{w_1}{w_0}\cdot\frac{\pi_s\,p_s(x\mid 1)}{(1-\pi_s)\,p_s(x\mid 0)}
  =\frac{w_1}{w_0}\cdot\frac{P_s(Y=1\mid X=x)}{P_s(Y=0\mid X=x)} .
\]
By Assumption~\ref{ass:srccal} the last ratio is $f(x)/(1-f(x))$, so
$P_t(Y=1\mid X=x)=T(f(x))$.  Since $T\circ f$ is $\sigma(f)$-measurable,
\[
  g_t(v)=\mathbb{E}_{P_t}\big[\,P_t(Y=1\mid X)\,\big|\,f(X)=v\,\big]=T(v),
\]
hence $\delta_{\mathrm{con}}=0$.
\end{proof}

Three consequences are worth stating separately.

\begin{remark}[What the decomposition buys]
$D_{\mathrm{label}}$ is a functional of $(\pi_s,\pi_t,\mu_t)$ alone.  It is
therefore removable with \emph{zero} target labels whenever $\pi_t$ is
identified from unlabeled data.  $D_{\mathrm{concept}}$ is by
Theorem~\ref{thm:decomp}(c) precisely what survives the best possible prior
correction, and no unlabeled procedure can estimate it: two targets with the
same $\mu_t$ can have any $g_t$ whatsoever.
\end{remark}

\begin{remark}[The interaction term is a warning, not a nuisance]\label{rem:R}
$R<0$ means the two shifts partially cancel.  A site in that regime looks
better calibrated than either component alone implies, and applying the free
correction there can \emph{increase} $\mathrm{CE}_1$.  Because $R$ is estimable
from the same labeled sample as $D_{\mathrm{concept}}$, the regime is
detectable rather than merely possible, and \S\ref{sec:results} reports it.
\end{remark}

\begin{remark}[Why $L^2$ is the primary endpoint]
Only the $L^2$ split is an identity; the familiar $L^1$ ECE obeys a triangle
inequality whose slack is itself shift-dependent.  Reporting $L^1$ alone would
make the components look non-additive for reasons having nothing to do with the
data.
\end{remark}

\subsection{Identification of the decomposition}
\label{sec:ident}

Estimating $D_{\mathrm{label}}$ without target labels requires $\pi_t$.  Let
$q_s(\cdot\mid y)$ be the law of $f(X)$ under $P_s(\cdot\mid y)$ and $q_t$ the
law of $f(X)$ under $P_t$ --- the latter observable from unlabeled target data.
Classical black-box shift estimation (BBSE) assumes \emph{pure} label shift,
under which
\begin{equation}\label{eq:bbse}
  q_t \;=\; \textstyle\sum_y \pi_t(y)\, q_s(\cdot\mid y).
\end{equation}
The premise of this paper is that \eqref{eq:bbse} fails across countries.  We
therefore work with the strictly weaker model
\begin{equation}\label{eq:relaxed}
  q_t \;=\; \textstyle\sum_y \pi_t(y)\, q_s(\cdot\mid y) \;+\; \eta,
  \qquad \eta(\mathcal{V})=0,\quad \eta\in\mathcal{E},
\end{equation}
where $\mathcal{E}$ is a declared class of admissible concept-shift
perturbations.  Write
$V_0=\operatorname{span}\{q_s(\cdot\mid y)-q_s(\cdot\mid y')\}$ for the space
of differences reachable by changing the prior, and
$\mathcal{D}(\mathcal{E})=\{\eta-\eta':\eta,\eta'\in\mathcal{E}\}$.

\begin{theorem}[Identification]\label{thm:ident}
The pair $(\pi_t,\eta)$ in \eqref{eq:relaxed} is uniquely identified from
$\big(q_t,\{q_s(\cdot\mid y)\}_y\big)$ if and only if
\[
  \{q_s(\cdot\mid y)\}_{y}\ \text{are linearly independent}
  \qquad\text{and}\qquad
  V_0\cap\mathcal{D}(\mathcal{E})=\{0\}.
\]
\end{theorem}

\begin{proof}
Two admissible pairs $(\pi,\eta)$ and $(\pi',\eta')$ yield the same $q_t$
exactly when $\sum_y(\pi-\pi')(y)\,q_s(\cdot\mid y)=\eta'-\eta$.  The left side
is a generic element of $V_0$, because $\pi,\pi'$ are probability vectors and
so $\pi-\pi'$ ranges over the sum-zero vectors; the right side is a generic
element of $\mathcal{D}(\mathcal{E})$.  Non-uniqueness is therefore equivalent
to the existence of a nonzero common element, i.e.\ to
$V_0\cap\mathcal{D}(\mathcal{E})\neq\{0\}$, together with the possibility
$\pi\neq\pi'$ inducing the zero element, which is excluded precisely by linear
independence.
\end{proof}

\begin{corollary}[BBSE as the degenerate case]
$\mathcal{E}=\{0\}$ gives $\mathcal{D}(\mathcal{E})=\{0\}$, so identification
reduces to linear independence of $\{q_s(\cdot\mid y)\}$ --- the standard BBSE
condition.
\end{corollary}

\begin{corollary}[Anchor relaxation]\label{cor:anchor}
Let $\mathcal{A}$ be a measurable \emph{anchor region} of score space and
$\mathcal{E}_{\mathcal{A}}=\{\eta:\eta|_{\mathcal{A}}=0\}$, i.e.\ concept shift
is assumed absent on $\mathcal{A}$ but unrestricted elsewhere.  Then
$(\pi_t,\eta)$ is identified if and only if the restricted family
$\{q_s(\cdot\mid y)|_{\mathcal{A}}\}_y$ is linearly independent.  This is a
\emph{rank condition on source data alone} and is strictly weaker than BBSE's
requirement that concept shift be absent everywhere.
\end{corollary}

\begin{corollary}[Norm-bounded shift is not point identified]\label{cor:partial}
If $\mathcal{E}=\{\eta:\|\eta\|_{1}\le\gamma\}$ for $\gamma>0$, then
$\mathcal{D}(\mathcal{E})$ contains a ball, so
$V_0\cap\mathcal{D}(\mathcal{E})\neq\{0\}$ whenever $V_0\neq\{0\}$ and the
model is \emph{not} point identified.  The sharp identified set for $\pi_t$ is
\[
  \Pi(\gamma)=\Big\{\pi_s^\top\! \operatorname{diag}(e_k)\,w \;:\;
     w\ge 0,\ \pi_s^\top w=1,\ \|q_t-Aw\|_1\le\gamma \Big\},
\]
an interval computable by two linear programs.
\end{corollary}

Corollary~\ref{cor:partial} is why the empirical sections report intervals and
a sensitivity sweep rather than a single corrected prevalence: $\gamma$ is an
assumption a deploying site must state, not a quantity it can estimate.  The
smallest $\gamma$ compatible with the data,
$\gamma_{\min}=\min_{w\ge0,\,\pi_s^\top w=1}\|q_t-Aw\|_1$, is a lower bound on
the true $\|\eta\|_1$ and is reported alongside; it is \emph{not} an estimate of
it, since BBSE chooses $w$ to make that residual small.

\begin{proposition}[An unlabeled test for concept shift]\label{prop:gof}
Discretise scores into $M>K$ bins, giving $A\in\mathbb{R}^{M\times K}$ with
$A[m,y]=P_s(\text{bin}=m,\,Y=y)$.  Under $H_0$ of pure label shift, $q_t$ lies
in the cone $\{Aw:w\ge0\}$, an $M-K$ dimensional restriction.  The pooled
Pearson statistic
$T=n_t\min_{w\ge0}\sum_m (\hat q_t[m]-(Aw)[m])^2/(Aw)[m]$
therefore has a testable null, and $H_0$ can be rejected using \emph{unlabeled
target data only}.
\end{proposition}

Proposition~\ref{prop:gof} is the most immediately deployable statement in this
paper: a hospital can evaluate it on its own unlabeled ECGs, before any chart
review, and learn whether a free prior correction will suffice or whether it
must budget for labels.  Two implementation points matter and are documented
with the software.  First, when $A$ is itself estimated, the variance
correction is \emph{per bin} --- the source contributes
$\sum_y w_y^2 A[m,y]/n_s$, weighted by the importance weights and far from
uniform across the score range.  Second, the asymptotic $\chi^2_{M-K}$
reference is anti-conservative at realistic sample sizes; we calibrate the null
by parametric bootstrap, which restores nominal size (measured
$0.035$--$0.065$ at nominal $0.05$) while retaining power $0.93$ against a
concept shift small enough to leave AUROC unchanged.

\subsection{Theorem 2: how many labeled target cases a site needs}

Fix a calibration budget $\varepsilon$ and a confidence level $1-\alpha$.  Define
\emph{calibration coverage} as
\[
  \Pr_{S\sim P_t^{\otimes n}}\Big[\;\mathrm{CE}_1(\hat h_S\circ f;P_t)\le\varepsilon\;\Big]\;\ge\;1-\alpha ,
\]
the probability being over the draw of the calibration set $S$.  This is
deliberately not a statement about the expected calibration error: a procedure
whose average is acceptable but which fails badly at one site in five is not
deployable, and an expectation hides exactly that.

Let $\Gamma=\mathrm{CE}_1(T\circ f;P_t)$ be the \emph{residual} concept shift
left after the free prior correction, and let $Z=\operatorname{logit}V$,
$p=\sigma(Z/\tau)$ under $\mu_t$.

\begin{theorem}[Sample complexity]\label{thm:samplesize}
\begin{enumerate}[label=(\alph*)]
\item \textbf{(Sufficiency.)}  For the hybrid estimator --- prior correction
followed by a one-parameter temperature fitted on $n$ labeled target cases ---
if $\Gamma<\varepsilon$ then calibration coverage holds as soon as
\[
  n\;\ge\;\frac{z_{1-\alpha/2}^2\,L^2}{I\,(\varepsilon-\Gamma)^2},
  \qquad
  L=\frac{\mathbb{E}[p(1-p)|Z|]}{\tau^{2}},\quad
  I=\frac{\mathbb{E}[p(1-p)Z^{2}]}{\tau^{4}} .
\]
\item \textbf{(A computable corollary.)}  By Cauchy--Schwarz $L^2/I\le\mathbb{E}[p(1-p)]$, so
\[
  n^\star\;\le\;\frac{z_{1-\alpha/2}^2\ \mathbb{E}[p(1-p)]}{(\varepsilon-\Gamma)^2}:
\]
\emph{the label bill is set by the model's mean predictive variance on the
target site, divided by the squared budget.}  Every quantity on the right is
computable from unlabeled target scores, so a site can price the work before
doing it.
\item \textbf{(Necessity, and a threshold.)}  If $\Gamma\ge\varepsilon$, any
procedure whatsoever requires
$n\ge(1-2\alpha)/(8\varepsilon^{2})$.
If $\Gamma<\varepsilon$, no such lower bound exists: $n^\star=0$ and the free
correction alone meets the budget.
\end{enumerate}
\end{theorem}

\begin{proof}
(a) Let $\tau^\star$ minimise the population log-loss over the temperature
family applied after $T$; $\tau^\star=1$ is feasible, so
$\mathrm{CE}_1(h_{\tau^\star}\circ f)\le\Gamma$.  For any $\tau$,
$|\sigma(z/\tau)-\sigma(z/\tau^\star)|\le \sup |\partial_\tau \sigma(z/\tau)|\,
|\tau-\tau^\star|$ and
$\partial_\tau\sigma(z/\tau)=-\sigma(1-\sigma)z/\tau^2$, so taking expectations
gives $\mathrm{CE}_1(h_{\hat\tau})\le\Gamma+L|\hat\tau-\tau^\star|$.  The
maximum-likelihood $\hat\tau$ satisfies
$\sqrt{n}(\hat\tau-\tau^\star)\Rightarrow\mathcal{N}(0,I^{-1})$, so
$\Pr[L|\hat\tau-\tau^\star|\le\varepsilon-\Gamma]\ge1-\alpha$ once
$\sqrt{nI}\,(\varepsilon-\Gamma)/L\ge z_{1-\alpha/2}$, which is the display.

(b) $\big(\mathbb{E}[p(1-p)|Z|]\big)^2\le\mathbb{E}[p(1-p)]\cdot\mathbb{E}[p(1-p)Z^2]$,
so $L^2/I\le\mathbb{E}[p(1-p)]$.

(c) Two-point (Le Cam) argument.  Construct $P_0,P_1$ with the \emph{same}
$\mathcal{X}$-marginal --- so no unlabeled procedure can distinguish them ---
and conditionals $g_0,g_1$ with $\|g_0-g_1\|_{L^1(\mu_t)}=2\varepsilon$.  Any
$\hat h$ with $\mathrm{CE}_1\le\varepsilon$ under both would contradict the
separation, so a procedure failing with probability $<\alpha$ under both yields
a test of $P_0$ versus $P_1$ with error $<\alpha$, which forces
$\mathrm{TV}(P_0^{\otimes n},P_1^{\otimes n})\ge1-2\alpha$.  Pinsker and
tensorisation give
$\mathrm{TV}(P_0^{\otimes n},P_1^{\otimes n})^2\le \tfrac{n}{2}\mathrm{KL}(P_0\|P_1)\le 2n\varepsilon^2$,
whence $n\ge(1-2\alpha)^2/(8\varepsilon^2)\ge(1-2\alpha)/(8\varepsilon^2)$ for
$\alpha\le 1/2$.  For the threshold: both hypotheses must lie in the class of
admissible concept shifts, a set of $L^1$ radius $\Gamma$ about $T\circ f$; two
of its points are at most $2\Gamma$ apart, so a separation of $2\varepsilon$ is
constructible only if $\Gamma\ge\varepsilon$.
\end{proof}

\begin{remark}[Where the labels go]
The content of (a) is the position of $\Gamma$ in the denominator.  Plain
temperature scaling spends its labels undoing a prior shift whose size is
already known from unlabeled data, and only then, with whatever precision
remains, on the concept shift.  Warm-starting at $T$ means every label is spent
on the part that genuinely requires them, so the requirement scales with the
\emph{residual} shift rather than with the total.  In \S\ref{sec:results} plain
temperature scaling with 300 labeled cases leaves calibration worse than
shipping the model unchanged, for the structural reason that a single
temperature cannot reverse the sign of miscalibration across the score range,
which a prior shift demands.
\end{remark}

\subsection{Theorem 3: prediction sets that stay valid}

A probability is one guarantee; a decision is another.  Let $s(x,y)$ be a
nonconformity score and $C_\alpha(x)=\{y: s(x,y)\le\hat q\}$.  Let
$P_{s\to t}$ denote the label-shift-corrected source, i.e.\ the law with
conditionals $P_s(x\mid y)$ and prior $\pi_t$.

\begin{assumption}[$\gamma$-bounded concept shift]\label{ass:gamma}
$\mathrm{TV}\big(P_t,\,P_{s\to t}\big)\le\gamma$.
\end{assumption}

\begin{theorem}[Conformal transfer]\label{thm:conformal}
\begin{enumerate}[label=(\alph*)]
\item \textbf{(Unlabeled.)}  Weighted split conformal on $m$ \emph{source}
calibration points with weights $w(y)=\pi_t(y)/\pi_s(y)$, run at level $\alpha$,
has target coverage $\ge 1-\alpha-\gamma$.  Under pure label shift
($\gamma=0$) it is exactly valid and requires no target labels.
\item \textbf{(Hybrid, Theorem 3.)}  Place total mass $\lambda$ on $n$ target
calibration points and $1-\lambda$ on the reweighted source points, and run at
the inflated level
\[
  \alpha' \;=\; \alpha-(1-\lambda)\gamma .
\]
If $\alpha'>0$ then target coverage is $\ge1-\alpha-O\big((n+m)^{-1}\big)$.
\item \textbf{(Efficiency.)}  The hybrid quantile has variance
$O\big(1/(\lambda^2 n+(1-\lambda)^2 m_{\mathrm{eff}})\big)$ against $O(1/n)$ for
target-only conformal, where $m_{\mathrm{eff}}$ is the weighted effective
source size.
\end{enumerate}
\end{theorem}

\begin{proof}
(a) Under pure label shift the reweighted source sample is exchangeable with a
sample from $P_{s\to t}$, so weighted split conformal
\citep{tibshirani2019conformal} gives coverage $\ge1-\alpha$ under
$P_{s\to t}$.  Coverage is the expectation of a $[0,1]$-valued indicator, so it
changes by at most $\mathrm{TV}(P_t,P_{s\to t})\le\gamma$ when the measure is
changed to $P_t$.

(b) Realise the weighting as a mixture: each calibration point is drawn from
$P_t$ with probability $\lambda$ and from $P_{s\to t}$ otherwise, so
calibration points are i.i.d.\ from $Q_\lambda=\lambda P_t+(1-\lambda)P_{s\to t}$
and the weighted empirical quantile is the quantile of $Q_\lambda$ up to the
usual $O((n+m)^{-1})$ finite-sample correction.  Since
$\mathrm{TV}(P_t,Q_\lambda)\le(1-\lambda)\gamma$, part (a)'s argument gives
coverage $\ge 1-\alpha'-(1-\lambda)\gamma=1-\alpha$.

(c) The quantile is an order statistic of a weighted empirical distribution;
its asymptotic variance is inversely proportional to the squared-weight-summed
sample size, giving the stated rate.
\end{proof}

\begin{remark}[Reading the trade-off]
$\lambda=1$ recovers target-only conformal, exactly valid but with a quantile
estimated from whatever handful of target labels exist; $\lambda=0$ recovers
the unlabeled reweighted-source procedure, stable but off by $\gamma$.  The
inflation $\alpha'$ makes the price of leaning on the source explicit and
payable in advance rather than discovered in deployment.  Note also the hard
floor no method evades: split conformal needs $n\ge1/\alpha-1$ target labels to
produce a non-trivial set at all --- nine cases at the 90\% level, ninety-nine
at 99\%.
\end{remark}

\begin{remark}[$\gamma$ is an assumption]
As with Corollary~\ref{cor:partial}, $\gamma$ cannot be estimated from the data
it is used on.  We therefore report a sweep over $\gamma$ rather than a single
coverage number, and a deployment dossier should contain that sweep.
\end{remark}
